\section{Make a wrong translation fail on purpose}
\label{sec:tiny-fault-lab}

A passing test suite tells us that the tested programs met the tested conditions.
It does not tell us whether those conditions would notice the mistake we fear.
We can ask a more concrete question: if multiplication were translated as addition, which test would fail?
This is not a proposal to damage the working compiler.
It is a controlled experiment with a separate, deliberately faulty component.

The program \code{fault\_lab.c} uses the correct compiler to emit stack code, then changes every \code{MUL8} instruction to \code{ADD8}.
It leaves each instruction's argument field alone.
It also leaves the source tree alone.
The fault lives only in this laboratory program; \code{compile\_expr} in \code{stack.c} remains unchanged.
Keeping the specimen separate makes it harder to confuse a demonstration of a bug with the implementation the reader should use.

Listing~\ref{lst:tiny-mutant} shows the faulty wrapper.
Both arithmetic instructions require two stack values and produce one.
Both require a zero argument field.
Replacing one with the other therefore preserves the structural properties checked by our verifier.
The generated program remains executable by the defined stack machine, but it generally computes the wrong value.

\begin{listing}
  \inputminted[firstline=7,lastline=15]{c}{example-code/ch01/fault_lab.c}
  \caption{This isolated wrapper deliberately changes multiplication into addition.}
  \label{lst:tiny-mutant}
\end{listing}

The specimen illustrates a narrow form of \term{mutation testing}: introduce a controlled change and check whether the tests detect its consequence.
One mutation does not model every compiler failure.
It checks a particular weakness in the evidence.
Here the question is whether we compare arithmetic meaning, rather than merely accepting structurally legal code.
The distinction fits the broader separation between compiler assurance tasks discussed by Safonov~\cite{safonov-2010-trustworthy}.
The specimen and its reduction procedure here are original, small experiments, not a survey of that subject.

\subsection{Choose a witness with a known answer}
\label{sec:tiny-fault-witness}

Use this source expression:
\[
  ((x+7)\times(3+5))+11.
\]
Its tree has nine nodes.
It extends the earlier seven-node tree with a final addition and the literal $11$.
At input $9$, the correct calculation is
\[
  ((9+7)\times(3+5))+11=16\times8+11=139.
\]
Every intermediate result fits in $V$, so no reduction changes these numbers.
The faulty target instead adds $16$ and $8$, then adds $11$, producing $35$.
We have a concrete disagreement: the expected value is $139$, and the target value is $35$.

A \term{counterexample} to a universal preservation claim is one allowed source program and input for which that claim fails.
The word does not require a large random test or an exotic boundary value.
Our ordinary input $9$ is enough.
The source is valid, its operations are total, and the target passes structural verification.
None of those facts can remove the disagreement.

It is useful to freeze the input while investigating this failure.
If a reducer changes the program and the input together, it may keep finding failures without showing which change simplified the explanation.
Our first reduction procedure keeps the witness input at $9$.
Later, once the program is smaller, we can ask whether the input matters at all.
These are two separate simplification decisions.

\section{Define which failure a reducer must retain}
\label{sec:tiny-reduction-predicate}

A \term{test-case reducer} searches for a smaller input that still exposes a chosen failure.
It needs a precise predicate for that failure.
For this laboratory, let $M(e,a)$ hold when all of the following are true:
\begin{enumerate}
  \item The faulty compiler successfully emits the whole program for $e$.
  \item The emitted code passes structural verification.
  \item The executor successfully runs it at input $a$.
  \item The returned value differs from $\llbracket e\rrbracket_a$.
\end{enumerate}
We test source validity by construction: every candidate is made from valid tree pieces or valid literal leaves.
The tree remains finite and non-sharing.
The predicate therefore describes wrong-code behavior on a valid, executable case.
It does not count a compilation failure or a rejected target as the same failure.

Listing~\ref{lst:tiny-mismatch} implements the predicate.
Returning false after a compilation failure does not claim that the compiler is correct.
It says that this candidate does not satisfy the particular predicate we are preserving.
Different investigations need different predicates.
An investigation of a compiler crash should retain that crash, not quietly substitute a wrong numeric result for it.

\begin{listing}
  \inputminted[firstline=17,lastline=27]{c}{example-code/ch01/fault_lab.c}
  \caption{The predicate retains successful executions with disagreeing values.}
  \label{lst:tiny-mismatch}
\end{listing}

This distinction prevents a common reduction mistake.
Suppose deleting part of a source program introduces an undeclared name.
If any nonzero tool status counts as interesting, the reducer may preserve a syntax error instead of the original miscompilation.
The final file can then be smaller but useless for explaining the bug.
Our structured candidates avoid syntax errors, and our predicate rejects changes that no longer reach a wrong-code execution.

The predicate also relies on a reference.
Here that reference is the previously defined evaluator, with arithmetic checked independently on literal pairs and small worked examples.
If the evaluator were wrong, a disagreement would need further investigation rather than an automatic declaration that the compiler was at fault.
Reduction makes the evidence easier to inspect; it does not certify the reference implementation.

\section{Reduce a tree without breaking its ownership}
\label{sec:tiny-greedy-reducer}

The reducer considers four kinds of replacement at each node.
It can replace a binary node with its left child, replace it with its right child, or replace any node with the literal $0$ or $1$.
It tries these choices in that order.
After an accepted replacement, it restarts its search at the root.
If no replacement at the current node is accepted, it searches its children from left to right.

This is a \term{greedy algorithm}: it keeps the first acceptable local improvement rather than examining every possible smaller tree before deciding.
It is not an algorithm for finding the smallest possible counterexample in an unrestricted language.
The candidate choices and their order matter.
They define which smaller cases this search can reach.

We also need a measure that decreases whenever a replacement is accepted.
Node count alone is not sufficient because replacing the input leaf with a literal keeps the node count unchanged.
Let $N(e)$ be the number of reachable tree nodes.
Let $W(e)$ be the sum of the weights of its leaves, assigning weight $n$ to literal $n$ and weight $256$ to the input leaf.
Compare the pair $(N(e),W(e))$ lexicographically: prefer fewer nodes, then a smaller weight sum when the node counts agree.

Every accepted replacement must decrease this pair and retain $M(e,9)$.
Replacing a binary node by a child decreases node count.
Replacing an input by $0$ or $1$ decreases weight without changing node count.
Replacing literal $3$ by literal $1$ decreases weight as well.
Replacing literal $0$ by literal $1$ is not a decrease and is not even tested for preservation.
Both components are nonnegative integers, so a strictly decreasing sequence of these pairs cannot continue forever.

Algorithm~\ref{alg:tiny-greedy-reduce} describes the outer search.
The candidate enumeration restarts after each successful replacement because the tree and its useful sites have changed.
The algorithm terminates when none of its permitted single replacements both reduces the measure and retains the predicate.
That stopping condition is narrower than global minimality.

\begin{algorithm}
  \caption{Reduce a failing tree using a fixed set of local replacements.}
  \label{alg:tiny-greedy-reduce}
  \begin{pseudocode}
    \Require A valid finite tree $e$ with $M(e,a)$; fixed input $a$
    \Ensure Retain $M(e,a)$; no permitted single improving replacement remains
    \Repeat
      \State $\mathit{changed}\gets\mathit{false}$
      \ForAll{sites in preorder, trying child replacements then literals $0,1$}
        \State form a candidate tree $e'$
        \If{$(N(e'),W(e'))<(N(e),W(e))$ and $M(e',a)$}
          \State $e\gets e'$
          \State $\mathit{changed}\gets\mathit{true}$
          \State \Break
        \EndIf
      \EndFor
    \Until{not $\mathit{changed}$}
  \end{pseudocode}
\end{algorithm}

The implementation changes one node's stored structure temporarily.
It saves the original structure, installs the proposed replacement, and evaluates the predicate on the entire root.
If the candidate fails, it restores the saved structure.
Predicate evaluation does not fold or otherwise mutate the tree.
Consequently, restoring the one changed structure is enough to restore the previous reachable tree.

Replacing a node with a child's structure follows the caller-owned storage rule established earlier.
The child's storage remains alive.
Disconnected nodes are not freed during reduction.
The original laboratory tree uses stack storage that outlives every predicate call, and no candidate introduces a cycle or a shared reachable node.
This is not a general-purpose reducer for arbitrary graph ownership schemes.
Its storage reasoning depends on precisely the same non-sharing contract as the earlier folder.

For the nine-node case, the C counters used by the measure are comfortably within the minimum range of unsigned integers.
An unbounded reduction library would need a separate argument or checked arithmetic for those counters.
The executable specimen is a bounded laboratory, not an unchecked promise about trees of arbitrary size.

\section{Work through every accepted reduction}
\label{sec:tiny-reduction-walkthrough}

The first candidate replaces the root addition with its left child.
The trailing $11$ disappears.
The new expression is the earlier product, which still disagrees at input $9$: the correct value is $128$, and the faulty value is $24$.
This step reduces node count from nine to seven.

At the product root, neither child alone retains the failure.
Both children contain only addition and leaves, so the faulty wrapper changes no instruction in their translations.
Replacing the root with $0$ or $1$ also removes the fault.
Those candidates are rejected.
The search moves into the left subtree instead.

Replacing $x+7$ with its left child gives $x\times(3+5)$.
At the fixed input, the correct value is $9\times8=72$ and the faulty value is $9+8=17$.
The candidate is accepted.
Replacing the input leaf with literal $0$ then gives $0\times(3+5)$.
The correct result is $0$; the faulty result is $8$.
Node count stays five in this second step, but leaf weight decreases.

The right subtree can next be replaced by its left child, literal $3$.
The result $0\times3$ computes $0$ in the source and $3$ in the faulty target.
Replacing $3$ with $0$ would erase the disagreement: both $0\times0$ and $0+0$ produce zero.
That candidate is rejected.
Replacing $3$ with $1$ retains the disagreement and decreases weight, so it is accepted.

Table~\ref{tab:tiny-reduction-steps} records the full accepted path.
The emitted diagnostic program prints the same node counts, input, and value pairs.
We can therefore check the hand reasoning against an executable trace of the search rather than merely asserting that reduction succeeded.

\begin{table}
  \centering
  \caption{Five accepted replacements retain the wrong-code failure at input 9.}
  \label{tab:tiny-reduction-steps}
  \begin{tabular}{lrrr}
    \toprule
    Source expression & Nodes & Correct & Faulty \\
    \midrule
    $((x+7)\times(3+5))+11$ & 9 & 139 & 35 \\
    $(x+7)\times(3+5)$ & 7 & 128 & 24 \\
    $x\times(3+5)$ & 5 & 72 & 17 \\
    $0\times(3+5)$ & 5 & 0 & 8 \\
    $0\times3$ & 3 & 0 & 3 \\
    $0\times1$ & 3 & 0 & 1 \\
    \bottomrule
  \end{tabular}
\end{table}

The final expression no longer contains the input.
It disagrees for every allowed input, not merely for the original witness.
Its correct translation is a push of zero, a push of one, and multiplication.
The faulty translation ends with addition instead.
Both sequences pass the verifier, both reach a stack of depth two before the final instruction, and both finish with one value.
The only relevant difference is the meaning of that final operation.

Figure~\ref{fig:tiny-fault-diagnosis} places the two conclusions side by side.
A structural acceptance is correct: there is no malformed instruction or missing operand.
A semantic acceptance would be wrong: zero and one are distinct results in $V$.
Adding more checks of the same stack-depth property cannot discover that difference.
The evidence must cross the source-to-target semantic boundary.

\begin{figure}
  \centering
  \begin{tikzpicture}[node distance=13mm and 17mm]
    \node[basic block] (source) {$0\times1$};
    \node[basic block,below left=of source] (good) {PUSH 0; PUSH 1\\MUL8};
    \node[basic block,below right=of source] (bad) {PUSH 0; PUSH 1\\ADD8};
    \node[basic block,below=of good] (zero) {valid; result 0};
    \node[basic block,below=of bad] (one) {valid; result 1};
    \draw[cfg edge] (source) -- node[above left,font=\small] {correct} (good);
    \draw[cfg edge] (source) -- node[above right,font=\small] {faulty} (bad);
    \draw[cfg edge] (good) -- (zero);
    \draw[cfg edge] (bad) -- (one);
  \end{tikzpicture}
  \caption{The verifier accepts both programs; only one preserves the source value.}
  \label{fig:tiny-fault-diagnosis}
\end{figure}

\subsection{Say exactly how small the result is}
\label{sec:tiny-reducer-minimality}

For this particular mutant, a literal or an input node cannot fail.
Its compiled code contains no multiplication instruction for the wrapper to change.
A source tree with an internal binary node has at least three nodes.
The reduced counterexample has exactly three nodes.
It is therefore smallest by node count for this mutant within our source grammar.

That conclusion uses a separate argument about the fault.
The greedy algorithm itself did not establish it simply by stopping.
For another fault, a different local replacement or a coordinated change at two sites could reveal a smaller failing program.
Local irreducibility means that the allowed single-step candidates do not improve the current case.
Global minimality means that no smaller failing case exists in the entire stated search space.
The words describe different claims and require different evidence.

Even a smallest node count need not give a unique explanation.
Several three-node multiplications expose this mutation.
Our tie-breaking rule prefers small literal weights, and the fixed candidate order leads to $0\times1$.
A reducer that tries literal $2$ first could stop at a different representative if its rules did not later reduce that literal.
The exact output belongs to the specified search procedure, not to an unspecified notion of the one true minimal test.

\section{Repair the contract boundary, then keep the test}
\label{sec:tiny-fault-regression}

The reduced program points directly to the broken local step.
The compiler's preservation proof requires a source multiplication to end in a target multiplication.
Replacing it with addition changes the value equation used in that induction case.
The stack-prefix shape remains correct, but its appended value is no longer the promised value.
The failure is in the arithmetic correspondence, not in the proof's storage reasoning.

In the working compiler there is nothing to repair: the injected fault is confined to the laboratory wrapper.
The program checks that the correct compiler runs the reduced case successfully and agrees with source evaluation for all inputs.
It also checks that the mutant disagrees for all inputs.
Both assertions matter to the demonstration.
Without the first, the laboratory could merely compare two defective paths.
Without the second, the specimen might stop demonstrating the intended fault after an unrelated edit.

A \term{regression test} preserves a previously exposed failure as an executable check against its return.
The small expression is a good candidate because its expected answer can be justified without the original larger tree.
It should complement, not replace, the earlier literal-pair tests.
Those tests already cover multiplication of zero by one and many nearby operand pairs.
The laboratory explains why that coverage detects a semantic error which the verifier correctly accepts as structurally safe.

Keep the pipeline configuration with the test.
The faulty wrapper changes emitted code, and the laboratory deliberately compiles the original source tree without invoking the folder.
If we first folded $0\times1$ into literal zero, no \code{MUL8} would be emitted.
The faulty wrapper would have nothing to change, and the disagreement would disappear.
That disappearance would not show that the faulty translation was repaired.
It would show that the selected execution path no longer reaches the broken operation.

This is why a useful failure record includes the source, input, enabled passes, target configuration, expected observation, and actual observation.
It also includes whether compilation and structural verification succeeded.
For this case all arithmetic and configuration choices fit in a small standalone C program.
In a larger compiler, the same discipline prevents a minimized file from losing the flags or pass ordering needed to reproduce its original failure.

\subsection{Know what reduction costs}
\label{sec:tiny-reduction-cost}

A predicate evaluation recompiles and executes a candidate and also evaluates its source tree.
With $N$ source nodes and enough capacity, each of those steps takes $O(N)$ time here.
Computing a candidate's size measure is another tree traversal.
The number of candidate attempts depends on the tree, the accepted steps, and the search order.
It would be misleading to call the whole reducer linear merely because one predicate call is linear.

If there are $R$ complete search sweeps over trees of at most $N$ nodes, each sweep considers at most a constant number of replacements at each node.
Using the simple full-tree predicate, this gives an $O(RN^2)$ upper bound for those sweeps, with fixed-width arithmetic and constant-capacity append treated as constant-time operations.
The bound does not state that every candidate takes the worst-case time.
It states which repeated work the straightforward implementation can perform.
The decreasing measure proves termination, but does not promise a small value of $R$ for an arbitrary starting tree.

For our nine-node example, the exact accepted path is more informative than an asymptotic slogan.
There are five accepted replacements and a final unsuccessful sweep.
We print only accepted states so that the demonstration stays readable.
A debugging tool could additionally log rejected candidates and their reason for rejection.
Those records should distinguish a nondecreasing measure, a lost mismatch, and a failure to compile or execute.
They are different explanations for why a proposed reduction was not retained.

The completed experiment gives us a useful division of labor.
The verifier guards target structure.
The source evaluator and target executor expose a semantic disagreement.
The reducer makes that disagreement small enough to inspect.
The proof identifies the local correspondence that was broken.
The regression test keeps the discovered case available after the investigation ends.
No one of these components supplies all the evidence that the others provide.
